\section{Discussion}
\label{sec:discussion}
\subsection{Contributions}
This paper furthers current understandings of the mechanisms connecting culture and status group stratification by bringing attention to the socio-demographic correlates of an important phenomenon.  The literature on culture stratification has been primarily focused on the question of the determinants of lifestyle consumption patterns, with a view to establishing whether there still exists a correlation between cultural choices and social position.  This one-sided concern has led to the neglect of an equally important question: that of the behavioral implications of familiarity and engagement with the arts.  In this respect current research has generally failed to consider the existence of stratification-linked differences in the extent to which cultural knowledge is differentially ``put to use'' by members of different social strata as a resource for interaction \citep{dimaggio87, deighton92, gainer95, benzecry_collins14}.  To answer to this question, I synthesized insights from classical theoretical sources in the culture and stratification literature with recent empirical and conceptual advances in our thinking of the relationship between culture and social networks.  I argued that one of the most important behavioral implications of familiarity with culture and the arts is its role as the type of content that is useful for the generation of conversational interaction; in short, \emph{culture talk}.  This is a phenomenon that deserves as much empirical and theoretical attention the ``patterns of cultural choice'' \citep{peterson83} that have traditional monopolized the attention of culture and stratification researchers.   

This paper contributes to the study of culture and stratification in several ways: analytically, I developed a framework that partially meets my self-imposed challenge by highlighting the theoretical and empirical importance of the culture talk phenomenon. I began with the Simmelian legacy on the importance of seemingly ``purposeless'' (non-instrumental and non-goal-oriented) conversational content for the emergence of ``sociability'' (or the ``play form'' of association) and connected it to the more contemporary stress by cultural stratification theorists on ``entertainment talk'' as the cement that holds status groups together as micro-relational communities of acquaintance and informal membership \citep{collins75, dimaggio87, gainer95, purcell07, lizardo06a, rivera12}. Empirically, I advance research in the growing literature linking taste, lifestyle consumption and stratification. I do this by examining the socio-demographic structuring of culture talk as a constitutive component of lifestyle-based stratification in contemporary societies. The results  show that, consistent with prior theory in cultural stratification, culture talk is stratified by the same markers of social position and status that previous research has shown also characterize audience segmentation in the corresponding cultural domains. More specifically, I show that culture talk in all domains---whether commercial or the fine arts---is stratified by education, income, age and gender, and to some extent urban location \emph{net} of self-reported engagement in those domains.  Those individuals endowed with larger stocks of economic (members of the middle and upper classes) and cultural capital (the highly educated) are more likely to routinely engage in culture talk, the young are more likely to engage in culture talk than the old,  women are more likely to engage in culture talk than men, and urbanites (in some domains) are more likely to engage in culture talk than their suburbanite/rural-dwellers counterparts.   

\subsection{Implications for the Study of the Culture/Network Link}
This paper also contributes to the emerging literature on the culture/network linkage \citep{fuhse09, pachucki_breiger10}.  This emerging line of theory and research  has coalesced around a critique of older styles of network research that followed Simmel in viewing structural form as separable from cultural content and mid-twentieth century British Social Anthropology in conceiving of cultural meaning as an epiphenomenon constituted by more ``concrete'' network relations and status positions  \citep{wellman88}.  Recent reconsiderations of the culture-network linkage have decisively moved beyond this received picture in demonstrating that culture and meaning are in fact constitutive of network ties and proposing a more balanced dynamic interplay between relational form and cultural content \citep{vaisey_lizardo10}.  Nevertheless, the relevance of this theoretical re-thinking for substantive research on the topic of cultural stratification is not yet clear (but see \citet{schultz_breiger10}).  This paper provides a concrete substantive application of these recent insights that clarifies the implications of this new literature for the issues that matter to cultural stratification researchers.   I propose that aesthetic consumption experiences are one of the key contents constitutive of important forms of network connectivity in contemporary societies \citep{dimaggio87}.  

This argument and the results reported in this paper have broader implications contemporary theory on the connection between culture, connectivity and structured inequality \citep{erickson08}.  In this respect, the type of routine deployment of cultural knowledge in concrete interaction that I have labeled as culture talk can be thought of as an important \emph{relational mechanism} \citep{tilly98, gross09}, linking inequalities in the propensity to use aesthetic content as a conversational resource to inequalities in relational outcomes associated to social connectivity and status-based boundaries. These long-standing inequalities in connectivity (both in terms of range and diversity) have been duly noted in the network literature \citep{lin99, erickson08}, but have yet to receive a satisfactory explanation.  Building on the results reported in this paper, it can be argued that different rates of engagement in culture talk provides a micro-interactional mechanism accounting for the dynamic persistence of cross-group inequalities on the major dimensions of social connectivity.

First, under the proposed account, individuals who have a larger propensity to engage in culture talk should also have an easier time during the relationship-formation process \citep{lazarsfeld_merton54}.  This is a plausible micro-mechanism that can help us to account for the existence stubborn status-linked inequalities social capital.  For instance, it is well-known that education has been \citep{marsden87} and continues to be \citep{lizardo06a, diprete_etal11} one of the best predictors of network size in the United States---and other national settings \citep{mollenhorst_etal11}---today.   The higher-educated are also more likely to engage in culture talk than the less-educated (within-levels of consumption), and thus they should also be more likely to routinely use this as an interactional resource to form new relationships \citep{fine77, dimaggio87}. The same can be said for those socio-demographic groups known to have smaller networks than average such as the less-educated elderly \citep{cornwell_etal08}.  Culture talk propensities decline with age for this group, consistent with the claim that rates of relationship formation underlie these differences.

Second, persons who belong to socio-demographic groups that exhibit a larger propensity to engage in culture talk should also tend to be embedded in status segregated networks.   Social ties that are partially constituted via culture talk retain this property because they are dynamically sustained via conversational meanings and contents access to which is also status-segregated \citep{collins75}.   In this respect, differential rates of culture talk across socio-demographic groups emerge as a plausible mechanism that accounts for the existence of status-segregated social networks even in the absence of any strictly conscious attempt on the part of individuals to select status-homogamous contacts.   

Third, it is well-established that the social-networks of Americans are characterized by high-rates of status-endogamy specially premised on race and education \citep{diprete_etal11}, and that marital homogamy premised on education and income has increased at the expense of marital endogamy premised on other axes of difference \citep{ kalmijn94, schwartz10}.  Persons who belong to high-education or high-income groups are also more likely to engage in culture talk than the less-educated (within-levels of consumption).  Following the same line of reasoning, they should also be expected to be more likely to routinely use this as an interactional resource to recurrently recreate status-segregated conversational rituals \citep{collins04, rivera12, nagel_etal11}.

Finally,  those individuals who have a more pronounced propensity to engage in culture talk should also tend to be more likely to bridge structural and ``cultural holes'' \citep{pachucki_breiger10, lizardo06a, lizardo11a}. This is a phenomenon that has been referred by \citet{schultz_breiger10}---riffing on \citet{granovetter73} and building on \citet{erickson96} and \citet{lizardo06a}---as \emph{the strength of weak culture}.  Status-linked inequalities in culture talk propensities can thus be proposed as a mechanism that explains the often noted fact that high-status individuals are more likely to be connected to a more diverse pool of others than low status individuals \citep{marsden87, erickson96, erickson08, lin_dumin86, lin99, goldstein_warren00, diprete_etal11}.   With persons of high education and/or high-income more likely to engage in culture talk than the less-educated, it stands to reason that they should also be able to---in some settings \citep{purcell07}---use this cultural expertise to make sporadic connections \emph{across} group and status boundaries.   

While the validity and scope of these hypotheses remains to be more fully explored, the results reported in this paper---along with convergent evidence from other sources \citep{lizardo06a, lizardo11a, lewis_etal09, rivera12, schultz_breiger10}---strongly point in this direction, suggesting exciting prospects for future research in this area.  In this manner, a focus on the culture talk phenomenon can serve to cross-fertilize research in the culture and stratification literature and the sociology of taste with research on inequalities in network access and social capital.  
